With the advancements in computational power, modern scientific programming practices, community-driven open-source and data-sharing initiatives, computational materials modelling is now stepping away from simply computing the properties for a small set of structures for ones with some prescribed compositions, to a paradigm in which high-throughput computations and statistical tools (broadly encompassing both PES samplings and ML/AI techniques) are applied jointly to holistically investigate the structure-energy-property relationships for a given material system. In particular, this is empowered by the developments and applications of ML and AI tools that help us accelerate the energy and property computations, enhance our understanding of the computational results in high dimensional space, and ultimately, to be used to suggest new research directions, specifically manifests itself through the proposal of new, stable material candidates with targeted properties.

These integrated approaches have impact significantly how atomistic modelling is driving the development of materials in the new era. For examples, it has enabled the establishment of holistic structure-energy-property relationships to understand how atomistic structures translate into the physical properties of materials. It promotes interdisciplinary researches, in which mathematical and statistical tools are developed in a material-specific context to help us derive agnostic, unbiased and physically informed materials representations and associated algorithms, to facilitate the establishment of these understandings, that are subsequently captured to build models to accelerate the sampling and predictions of materials properties. 

These advancements have played a critical role in enabling us to tackle much complex problems in materials physics and chemistry that traditional methods suffer, predominantly limited by the approximations for the corresponding theories were adopted, in order to convert these theories into practically workable computational tools. We have already seen promising developments in breaking these limitations across multiple fronts, such as in the investigation of electrocatalytic surfaces under explicit solvents,\cite{Qin2023CO2Electroreduction,Sakong2018ElectricDoubleLayer,Andrade2020ProtonTransfer} electro- and magneto-caloric effects for energy conversions,\cite{graf2021unified} phonon-mediated superconductivity beyond harmonic phonons,\cite{zappacosta2025nonlinear} to name just a few.

By providing a pedagogical review of the established methodologies, by balancing the mathematical formality, theoretical rigorous, and intuitive explanations, we hope that the present monograph can facilitate the computational materials practitioners to further: (a) enhance their understanding on the working principles, capabilities, and limitations behind the existing methodologies in a conceptually unified framework, (b) build competency in the applications of the appropriate methods for solving problems in a more physically informed way, and (c) equip themselves to capture the emerging capabilities of modern AI tools for tackling challenging problems in materials physics and chemistry.

Looking forward,  there are still plenty of open questions to be answered, which are providing the momentum for the continuous developments of  atomistic molecular and materials modellings. In the final section of this monograph, some of these open questions that we believe are most relevant to the topics being covered here  shall be raised. One could easily appreciate that, many of these questions are cross-cutting problems that, each one of them  correlate multiple topics that have been discussed here. 

\textit{\textbf{What are the roles of transitional methods, such as empirical atom-atom potentials and DFT, in materials modelling, with respect to the development of more generalised, broadly applicable AI models?}} Modern deep-learning potentials, for examples, multiple rounds of MP operations in GNN-based potentials, are usually computationally demanding, and require GPU accelerations to make them competitive with non-accelerated empirical atom-atom potentials.\cite{leimeroth2025machine} This makes classical forcefields still possess competitive advantages when it comes to atomistically simulating very large systems close to a mesoscopic scale. Most importantly, even with foundational energy models,\cite{batatia2023foundation} our pursuit of scientific discoveries meaning that more than often we will be dealing with problems that are out-of-distributions in the training set of the foundation models. For examples, in the studies of anharmonic solids,\cite{yang2026mace} fast-ion conductors,\cite{ding2025liquid} where the ions frequently visit lattice positions/chemical environments that are not the equilibrium positions presented in the training data, structure and dynamics of magnetic materials where foundational models for these systems are still at the emerging state,\cite{ho2026equivariant} as well as the modelling of complex interfacial phenomena in which the charge, spin and orbital interactions interwind that leads to unconventional physical properties. Systems such as these still require detailed and reliable electronic structure theory calculations, to facilitate the modelling of these systems and data generations and validations for training future AI models. 

Most importantly, the continuous developments of the electronic structure theory, which is heavily grounded by the fundamental physical principles, rather than data representations, still play a pivotal role in facilitating the understanding of exotic physical phenomena in material systems. At the same time, we are yet to discover more routes to incorporate these fundamental knowledge in materials physics and chemistry into future AI models. Hence, it is still very immature to think AI is taking over the traditional tools in atomistic modellings.  

\textit{\textbf{How to know if one have sampled enough?}} This is probably one of the most persistent and challenging problem to answer in all investigations that involve PES sampling. Certainly this is a very problem-specific question, meaning to a large extent, it depends on the nature of the problems that we are trying to solve. If one is trying to determining the free energies of high-entropy alloys, in which we know the material does not confine into a specific form of atomic arrangements, the convergence of the statistical sampling can be gauged by monitoring the evolution of the probability distributions of the energies of the local minima being sampled. This may be considered as a relatively \emph{loose criterion}. On the other hand, in the case such as organic crystal structure predictions,\cite{case2016convergence}  where there is a strong interest in locating the global minimum (a \emph{strong} criterion), then one may have to keep running the sampling until the same global minimum has been repetitively visited for multiple times. 

The generally lack of a robust check for the convergence also affects our capabilities in achieving predictive computations of certain materials properties through statistical sampling, such as the computation of lattice thermal conductivities in real space with Green-Kubo method [\cref{tab:properties_from_MD}].\cite{knoop2023ab} In materials where long wavelength $|\bm{q}|\to 0$ phonons (that needs to be modelled with an indefinitely large supercell) can be  strongly scattered  to suppress $\kappa_L$, simply using larger and larger supercells in the calculation may never converge $\kappa_L$ or converge it to a wrong value, and the Green-Kubo method itself does not necessarily provide us with a diagnostic that we are dealing with such a kind of materials.  

\textit{\textbf{Can generative modelling help us to sample better?}} It is therefore natural to ask, whether the development of generative models can help us to sample better. We are still yet to see more evidences that will allow us to answer this questions in a more objective way. Nevertheless, do not forget that, on the very fundamental level, in method such as the score-matching diffusion model,\cite{song2019generative} the generative process is still largely \emph{guided} by a forcefield\cite{cheng2024response} that are trained predominantly on the equilibrium distributions of stable materials. Whether this imposes a limit on our capability in exploring the chemistry and structural landscape of the materials, is probably worth further investigations. Answering this question, in the context of this monograph, would help us to unleash the greater promises of the generative models that we are yet to capitalise. 

\emph{\textbf{What are the future perspectives and promises behind the development of DFT and higher-level electronic structure methods for materials modelling?}} The 2022 ``DFT exchange'' paper,\cite{dftexchange} has gathered the leading theoreticians behind the developments of DFT and related theoretical chemistry methods to share their views behind the current status and future trends in DFT from a theoretical perspective. The present monograph is partially inspired by this unique publication, and trying to focus on a more practical side on the applications of DFT. In this aspect, some of the ongoing challenges in electronic structure theory that are relevant to the atomistic materials modellings include (but are not necessarily limited to) the followings:

(a) \emph{What are the perspectives on the development of more accurate exchange-correlation functionals?} Perhaps of the most interests to materials scientists would be the perspective of having a universally good functional to achieve reliable predictions on the electronic band gaps across a wide range of materials, by overcoming the self-interaction errors but at a computational cost that is comparable to GGA. Tackling this problem would certainly depending on the viewpoints of the DFT developers on whether one should achieve a universal functionals that works for all systems, or one should tailor-design them  for different categories of systems with their unique physics.\cite{mardirossian2016accurate} This is not a judgement that we tend to make or suggest here.

(b) \emph{What are the roles of the experimental data for fitting and benchmarking density functionals?} %As a computational materials scientist, a repetitive argument that I received during project discussions,  peer review of journal articles and grant proposals, is on how well the results from computational modelling agree with experimental results, if not, why model. One would then going on to argue that one should use experimental data, rather than computationally generated ones, to benchmark new theories, and train/validate new models. 
In the development and testing of new density functionals, the role of experimental data cannot be neglected, as the ultimate goal of atomistic modelling is for us to guide the experimental discoveries of materials. Nevertheless, one should always bear in mind that, (i) We should not expect a theory or model to be able to fully account all the aspects that may influence the experimental outcomes, as a complex model itself may not lead to valuable and explainable insights. (ii) The experimental data itself, even carefully curated, often need to be corrected, such as the thermal effects,\cite{otero2012benchmark} before they can be used to benchmark the theory, in which those effects are not accounted for. In addition, these experimental data should ideally be collated and corrected with a consistent set of settings, which  impose additional challenges on the experimental designs. 

Certainly different material systems and measurement conditions may require different correction scheme, and whether the contributions from different factors, such as defects, can always be isolated, will remain a highly material-dependent problem. Hence we expect the development on the aspect of experimental data may be lying something in the middle. The same argument  probably also applies to as to what extent future (foundational) AI models for materials can be trained on experimental data.

(c) \emph{AIMD only account for dynamics of nuclei under thermal excitation, how about electrons?} Thermal excitations not only affects the dynamics of nuclei, the electronic subsystems are also excited and can play a critical role in determining a variety of excited state behaviours in molecules and materials. However, such effect is not accounted in ground-state DFA. Recent developments on workable implementations of ensemble DFAs\cite{valone1980one} provides the community a workable route to tackle this problem, and we are yet to see more exciting and practical applications of such methods in molecular and materials modellings.

(d) \emph{What are the future directions of modelling the response behaviours under external stimuli?} Fundamentally, this question is linked to the excited state properties of materials, which have not been extensively discussed in the monograph. Amongst all the possible response properties of materials under external fields, materials scientists would probably argued that the modelling of magnetic responses in materials is properly one of the most interesting aspects to be looked at. The theoretical challenges behind the formulation of DFT in a magnetic field are still nevertheless enormous,\cite{DFT_exchange} so we are still to see where the future promises in this aspect are lying.

(e) \emph{How about the many-body methods?} As we have already demonstrated here, GW method has become increasingly important and popular in modelling materials properties in which many-body effect plays a critical role, such as the optical properties of 2D materials. A more accurate, faster and scalable GW methods, probably with the help of AI models,\cite{hou2024unsupervised,hou2025data} that are more integrated with computational workflow managements,\cite{ganose2025atomate2} would be important to make these methods more accessible to materials modelling practitioners. To make this happens would also largely depend on one final question raised as following.

(f) \emph{How can we better integrate atomistic  materials  modelling into the existing educational program of chemical and materials science?}  Successfully achieving this will be critical to drive the future developments of this field. A critical challenge is likely to be how the educators can convey and deliver the difficult concepts in physics, highly abstract ideas in mathematics and practical skills in computational programming in an approachable way, such that the students can slowly build up their confidences in mastering the many faceted fundamentals that underpins atomistic materials modellings, which will ultimately be translated into algorithms and research software to deliver meaningful and valuable outcomes. 

\textbf{\emph{What are the future promises in the application of atomistic modelling in understanding and predicting materials' properties?}} With the increase in computational power, the development of new materials theory and integrated workflow managements for high-throughput computation, together with better education and advocation of good atomistic modelling practices, we should see an increasing adaptation and integration of the existing and future tools in solving many-body problems in materials science beyond a mean-field or lower-order perturbative approaches. This will significantly wider the scope over the materials systems and problems we can tackle. Many of these had been encountered in the previous discussions, such as the thermal properties of highly anharmonic solids, superconductivities mediated by anharmonic phonons, and many others. Advances in this front will also enable us to curate more highly specialised and accurate dataset for training ML/AI models that cover a wider range of materials properties.


We now move on to raise some more specific questions related to the development of ML/AL in materials science.

\textbf{\emph{How should we treat the `zoo' of AI models for materials modelling? }} DFT practitioners are often said to be challenged by a `zoo'\cite{goerigk2017look} of exchange-correlation functionals that have been developed over many decades, each one being developed as some approximation to the complex quantum behaviours of electrons, that are relevant in the chemical systems (\emph{either broadly or specifically to some specific ones}). Their successes and pitfalls are identified through extensive benchmarks across a wide chemical space.

As one can see in \cref{tab:equivariant_GNNS} and \cref{sec:generative_AI}, the number of AI models developed for atomistic molecular and materials modelling is also rapidly increasing over the past few years, with many of them not only being developed by chemists and materials scientists, but also increasingly by mathematicians and computer scientists who are not traditionally working in this field. In response to this, there is an increasing focus in benchmarking the performance of different AI models against a consistent training and validation set to facilitate the community in the selection of best performing models for materials modelling,\cite{leimeroth2025machine,eastman2026accuracy,yu2024systematic,focassio2025performance,deng2025systematic,riebesell2025framework} which is definitely  much needed information. However, a more deeper scientific questions behind these benchmarking is that whether some numerical performance indicators (such as prediction errors), though simple to compare, may also lead to a similar problem of `\emph{choice by popularity}' with the use of B3LYP\cite{kruse2012standard} in DFT practices, as more users are converged into the `\emph{best performing}' models without the desire to properly understand why these models perform well. We envisage that more analyses, tools and frameworks that are formulated and developed to facilitate the community in understanding how subtle differences behind model designs can lead to the differences in model performances will definitely be very beneficial.

\textbf{\emph{Are foundational models for materials modelling truly `foundational'?}} Needless to say, the concept of foundational models are borrowed from the Large Language Models,\cite{vaswani2017attention} that lead to successful applications such as ChatGPT. It is foundational in a sense that it is trained on all the available text-based resources so it can infer to almost all range of available human knowledge. In a sense that it is a universal model. 

However, it is very obvious that in materials physics and chemistry, we do not yet have a universal theory (or an approximation of which), that is applicable to all application scenario for all existing and to-be-discovered materials. Hence, the word foundational model when applied to this specific scientific disciplinary may lead to some incorrect perceptions on where we are currently standing on.  At the time of writing this monograph, the problems that are common to all ML/AI models are still highly relevant to foundational models, being (i) the ability for the models to extrapolate is highly correlated with the data on which the models are being trained on,\cite{yang2026mace,fu2022forces} and (ii) the models' accuracies are also highly correlated with the underlying theory (\emph{e.g.} the exchange-correlation functionals\cite{shi2026influence}) that are adopted to generate the training data. Hence, one still needs to be cautious when using the foundational models, as one would approach the problem of choosing appropriate DFAs to use in a similar way. Nevertheless, foundational models have indeed provided us a faster alternative to DFT for pre-screening the material space to facilitate the subsequent, more accurate downstream  tasks in atomistic modellings.

\textbf{\emph{To what extent the trained AI models can help us discovering new knowledge in materials science at the atomistic level?}}
NN are often considered as another type of `black-box' method that maps an input to some desired output, whereas how and what has been learnt by the models remain largely under-explored. Models such as GNN\cite{duval2023hitchhiker} and attention layers,\cite{vaswani2017attention} in principle, are trying to encode atom-atom interactions into a simpler latent space that may facilitate the understanding of which interactions are important, for example, in determining materials stabilities. For examples, \citet{hafizi2026towards} analysed the attention matrices learnt by the NN model and tried to correlate them with the stabilities of organic crystal structures.
On the other hand, one may also argue to what extent the information encoded in these matrices are meaningful as they may be just  an ansatz corresponding to a plausible local minimum on a high dimensional landscape of the loss function. We believe that the continuous developments in the `\emph{Explainable AI}'\cite{minh2022explainable} will offer us with more opportunities in solving this problem. Another promising direction is behind the development of \emph{neural operators},\cite{azizzadenesheli2024neural} which aims at better incorporating existing physical knowledges, such as electrostatics, electron densities and wavefunctions, into the model design to enhance their capabilities in solving more physically relevant tasks, this could simultaneously help to make the model becoming more explainable. 

%\textbf{\emph{What are the future promises of the generative models?}}


\textbf{\emph{Will the AI models eventually have the capability to discover equivariance and the high-level (such as the graph structures) of materials themselves, rather being enforced as an inductive bias?}} The final questions that we believe is worth asking, is on the importance of enforcing symmetry in the design of AI models for materials modelling, which, in this monograph, we have covered extensively on how symmetry can be incorporated in both ML/AI models to enhance models' data efficiency and abilities to extrapolate. However, the counter argument behind this is that, in a truly `data-driven' paradigm, the model should be able to discover these symmetries and enforcing them, without explicit data augmentations,\cite{duval2023hitchhiker} which can be noted in some recent research being devoted to address such a problem.\cite{domina2026unconstrained,bigi2026pushing,langer2024probing} 

We believe that the ongoing theoretical advances across all fields related to atomistic materials simulations in the coming years will continuous to provide us with exciting discoveries related to the above questions. 